
% One-slide talk (O-slide): основная идея работы на одном слайде.
% Числа и рисунок воспроизводятся командой: python code/main.py
% Разместить файл в папке slides/ проекта.
\documentclass[aspectratio=169]{beamer}

\graphicspath{{../reports/figures/}}

\usepackage{cmap}
\usepackage[T2A]{fontenc}
\usepackage[utf8]{inputenc}
\usepackage[english,russian]{babel}
\usepackage{amssymb,amsfonts,amsmath}

\usecolortheme{beaver}
\usefonttheme[onlymath]{serif}
\setbeamertemplate{navigation symbols}{}
\setbeamersize{text margin left=0.35cm,text margin right=0.35cm}
\setbeamerfont{frametitle}{size=\normalsize}

\setbeamertemplate{footline}{%
  \hbox{\hspace{0.35cm}\tiny\color{gray}%
    Александр Ситников \quad МГУ имени М.\,В.~Ломоносова}%
  \vspace{2pt}%
}

% Сгенерировано code/main.py (этап tables) — не редактировать вручную.
\newcommand{\PilotAsset}{Si}
\newcommand{\PilotPeriod}{2025-10-01--2026-09-30}
\newcommand{\PilotDays}{254}
\newcommand{\PilotTrades}{21\,217}
\newcommand{\PilotTradesPerDay}{83}
\newcommand{\PilotMaskLow}{10}
\newcommand{\PilotMaskHigh}{90}
\newcommand{\PilotCbLow}{1{,}05}
\newcommand{\PilotHistCbLow}{1{,}04}
\newcommand{\PilotCbHigh}{3{,}78}
\newcommand{\PilotHistCbHigh}{1{,}21}
\newcommand{\PilotHistHigh}{1{,}29}
\newcommand{\PilotDiffLow}{-0{,}01}
\newcommand{\PilotDiffLowCI}{[-0{,}09;\ 0{,}08]}
\newcommand{\PilotDiffHigh}{-2{,}56}
\newcommand{\PilotDiffHighCI}{[-2{,}79;\ -2{,}37]}
\newcommand{\PilotVisibleLow}{75}
\newcommand{\PilotVisibleHigh}{8}


\begin{document}

\begin{frame}[t]{История снижает ошибку восстановления волатильности при дефиците сделок}

\vspace{-0.6em}
{\scriptsize
Пилот: опционы на фьючерс \PilotAsset{} (MOEX),
\PilotDays{} торговых дней, \PilotTrades{} сделок.
Скрытые сделки восстанавливаются по оставшимся наблюдениям.
}

\vspace{0.3em}
\begin{columns}[T,onlytextwidth]

  \begin{column}{0.67\textwidth}
    \centering
    \includegraphics[
      width=\linewidth,
      height=0.57\textheight,
      keepaspectratio
    ]{oslide_pilot_\PilotAsset}
  \end{column}

  \begin{column}{0.32\textwidth}
    \scriptsize

    \textbf{Результат пилота}
    \begin{itemize}\setlength{\itemsep}{0pt}

      \item Скрыто \PilotMaskHigh\,\% сделок:
      \PilotCbHigh{} $\rightarrow$
      \textbf{\PilotHistCbHigh{} п.п.}

      \item Только вчерашняя кривая:
      \PilotHistHigh{} п.п.

      \item Скрыто \PilotMaskLow\,\%:
      \PilotCbLow{} и \PilotHistCbLow{} п.п.
      (без значимого различия).

    \end{itemize}

    \vspace{0.15em}
    \textbf{Обозначения}
    \begin{itemize}\setlength{\itemsep}{0pt}

      \item $m$ --- доля скрытых сделок;

      \item $S_{t-1}$ --- вчерашняя
      расчётная кривая биржи;

      \item RMSE --- ошибка на скрытых
      сделках в п.п. IV;
      полосы --- 95\,\% доверительные
      интервалы.

    \end{itemize}
  \end{column}

\end{columns}

\vspace{-0.4em}
\begin{block}{}
  \scriptsize
  \textbf{Вывод:}
  в пилоте история существенно снижает ошибку
  при сильной разреженности данных.

  \textbf{Далее:}
  разработка метода, учитывающего историю,
  связи между контрактами и ограничения
  отсутствия статического арбитража.
\end{block}

\end{frame}
\end{document}
